\chapter{Étapes 4 et 5 — Fabriquer les surfaces et les contrôler}
\label{chap:maillage}

\begin{chapeau}
L'étape 4 transforme chaque os segmenté en une surface 3D fermée, un fichier STL. L'étape 5 mesure la sensibilité du
résultat au seuil choisi et écrit les rapports.
\end{chapeau}

\section{Remplir l'os}

Le STL doit décrire la surface extérieure de l'os, pas le canal médullaire. On remplit donc chaque os :
\begin{enumerate}
  \item une \textbf{fermeture morphologique} de rayon 1~mm (une dilatation suivie d'une érosion) bouche les petits
  pores de la corticale ;
  \item si l'os sort du champ du scanner, sa section au bord est d'abord fermée par un \textbf{couvercle} (la section
  remplie en 2D) : sinon le canal resterait ouvert et ne se remplirait pas ;
  \item un \textbf{remplissage des trous} en 3D remplit tout ce qui n'est pas relié à l'extérieur : canal et os
  spongieux.
\end{enumerate}

\section{Extraire la surface}

\subsection{Marching cubes}

Le masque rempli vaut 1 dans l'os et 0 dehors. On le floute légèrement ($\sigma = 0{,}3$~mm) pour obtenir un champ
$g$ qui passe en douceur de 1 à 0, puis on extrait la surface $g = 0{,}5$ par l'algorithme des \emph{marching cubes}
\cite{lorensen1987}, dans sa version rapide \emph{Flying Edges} \cite{schroeder2015}. Le principe : on parcourt les
petits cubes formés par 8 voxels voisins. Pour chaque cube, on regarde quels coins sont au-dessus de 0,5 ; parmi les
256 cas possibles, une table donne les triangles à placer dans le cube. Le sommet d'un triangle posé sur l'arête qui
relie deux coins $a$ et $b$ est placé par interpolation linéaire, au point
\begin{equation}
  \vb{x} = \vb{x}_a + t\,(\vb{x}_b - \vb{x}_a), \qquad t = \frac{0{,}5 - g_a}{g_b - g_a}.
\end{equation}
Le léger flou supprime l'effet d'escalier des voxels. Les sommets, calculés en indices de voxels, passent ensuite en
millimètres par la matrice $\mt{A}$ (\cref{eq:affine}).

\subsection{Lisser sans rétrécir}

Un lissage simple, qui rapproche chaque sommet de la moyenne de ses voisins, fait rétrécir le maillage. On utilise un
lissage « à fenêtre sinc » \cite{taubin1995} (20 itérations, bande passante 0,05) : c'est un filtre passe-bas sur la
position des sommets, qui enlève les petites bosses sans changer la taille de l'objet. C'est aussi le choix de 3D
Slicer.

\subsection{Alléger}

Au-delà de 150\,000 triangles (environ 7~Mo de STL), le maillage est simplifié par la méthode des \emph{quadriques}
\cite{garland1997}. Chaque sommet accumule une matrice $\mt{Q} = \sum_{p} \vb{p}\vb{p}^{\mathsf T}$, où
$\vb{p} = (a, b, c, d)$ décrit le plan $ax + by + cz + d = 0$ de chaque triangle voisin. L'erreur
$\vb{x}^{\mathsf T}\mt{Q}\,\vb{x}$ (en coordonnées homogènes) est la somme des carrés des distances de $\vb{x}$ à ces
plans. On fusionne en priorité les sommets dont la fusion coûte le moins. La simplification n'est gardée que si le
maillage reste étanche, et elle ne retire jamais plus de 75~\% des triangles.

\section{Placer la surface au vrai bord}
\label{sec:demi-hauteur}

\subsection{Le biais d'un seuil bas}

Le scanner floute les bords (\cref{sec:scanner}). Modélisons un bord d'os par une marche floutée : du côté de l'os
(position $s < 0$ le long de la normale) l'intensité vaut $C$ (corticale, environ 1300~HU), dehors elle vaut $B$
(tissus mous, environ 40~HU), et entre les deux
\begin{equation}
  I(s) = B + (C - B)\,\Phi\!\left(-\frac{s}{\sigma}\right),
  \label{eq:bord}
\end{equation}
où $\Phi$ est la fonction de répartition de la loi normale et $\sigma$ l'écart-type du flou. Le vrai bord est en
$s = 0$, là où l'intensité vaut la moitié entre $B$ et $C$. Un seuil $T$ place la surface là où $I(s) = T$, donc en
\begin{equation}
  \tcbhighmath{s_T = -\sigma\,\Phi^{-1}\!\left(\frac{T - B}{C - B}\right).}
  \label{eq:biais}
\end{equation}

\fig{profil_bord.pdf}{Un seuil bas (226~HU) place la surface trop à l'extérieur d'environ un écart-type du flou ; la
demi-hauteur retrouve le vrai bord.}{fig:profil-bord}

\begin{exemple}
Fantôme du test \fichier{test_chaine.py} : flou du scanner $\sigma_{\text{PSF}} = 0{,}5$~mm et lissage du logiciel
$0{,}35$~mm, soit $\sigma = \sqrt{0{,}5^2 + 0{,}35^2} = 0{,}61$~mm. Avec $T = 226$, $B = 40$ et $C = 1300$~HU :
$(T - B)/(C - B) = 0{,}148$, $\Phi^{-1}(0{,}148) = -1{,}05$, donc $s_T = +0{,}64$~mm. Le test mesure $+0{,}60$~mm : la
surface au seuil est trop grande de 0,6~mm, plus que la tolérance de 0,5~mm du projet.
\end{exemple}

\subsection{L'affinage à demi-hauteur}

Pour supprimer ce biais, chaque sommet est déplacé le long de sa normale jusqu'au point où l'intensité vaut la moitié
entre le pic de la corticale et le fond :
\begin{enumerate}
  \item on lit le profil d'intensité le long de la normale, de 2,5~mm à l'intérieur à 2~mm à l'extérieur, tous les
  0,1~mm (interpolation trilinéaire dans l'image) ;
  \item le \emph{pic} est le maximum du côté intérieur, le \emph{fond} le minimum après le pic ;
  \item le nouveau bord est le premier point, après le pic, où le profil passe sous $(\text{pic} + \text{fond})/2$
  (interpolation linéaire entre deux échantillons) ;
  \item le déplacement est ignoré si le contraste est faible (moins de 300~HU) ou s'il dépasse 1,5~mm (articulation
  serrée, profil douteux) ;
  \item les déplacements sont lissés sur le maillage (cinq passes de moyenne avec les voisins), pour ne pas créer de
  bruit d'un sommet à l'autre.
\end{enumerate}

\begin{exemple}
Sur le fantôme, l'écart passe de $+0{,}60$ à $+0{,}16$~mm ; le reste vient de la corticale très fine du fantôme
(1,5~mm), dont le pic est atténué par le flou. Sur le lévrier, les sommets reculent en moyenne de 0,21 à 0,32~mm
vers l'intérieur selon l'os.
\end{exemple}

\section{Vérifier que la surface est fermée}

Un fichier à imprimer doit décrire un solide fermé. Le logiciel vérifie :
\begin{itemize}
  \item l'\textbf{étanchéité} : chaque arête est partagée par exactement deux triangles ;
  \item l'\textbf{orientation} : les triangles voisins sont parcourus dans le même sens, et les normales sortent de
  l'os. Le volume, calculé par $V = \tfrac16 \sum_{\text{triangles}} \vb{a}\cdot(\vb{b}\times\vb{c})$, doit être
  positif ; sinon on retourne tous les triangles ;
  \item la \textbf{caractéristique d'Euler} $\chi = S - A + F$ (sommets, arêtes, faces) vaut 2 pour une surface fermée
  sans trou, comme une sphère. Un os seul doit donner $\chi = 2$ ; un groupe de $n$ os sans trou donne $2n$.
\end{itemize}

Les fichiers sont écrits au format \textbf{STL binaire} : 80 octets d'en-tête, le nombre de triangles, puis
50 octets par triangle (sa normale, ses trois sommets en réels de 4 octets et 2 octets libres). Le logiciel écrit aussi
\fichier{segmentation.seg.nrrd}, qui contient toutes les étiquettes avec leurs noms et leurs couleurs et s'ouvre dans
3D Slicer par-dessus la série DICOM, pour vérifier ou retoucher.

Enfin, l'étape 4 note les \textbf{bords du champ} touchés par chaque os (un plan et sa normale) : l'étape 6 s'en sert
pour savoir si une extrémité articulaire manque (\cref{sec:completude}).

\section{Étape 5 — La sensibilité au seuil}
\label{sec:sensibilite}

Le document de cadrage demandait de vérifier que le résultat ne dépend pas trop du seuil choisi (« variation du seuil
de ± 50~HU et mesure de l'impact »). Pour chaque os principal, le logiciel :
\begin{enumerate}
  \item re-segmente l'os à 176 puis à 276~HU, dans la même zone (l'os dilaté de 3~mm, sans les autres os) ;
  \item le remplit de la même façon et en extrait la surface ;
  \item mesure la distance de chaque sommet de cette nouvelle surface à la surface de référence.
\end{enumerate}
La distance est calculée avec une précision inférieure au voxel. Si $g$ est le masque de référence flouté (celui du
STL) et $\vb{p}$ un point proche de la surface $g = 0{,}5$, un développement au premier ordre donne
\begin{equation}
  g(\vb{p}) \approx 0{,}5 - d\,\norm{\nabla g(\vb{p})}
  \quad\Longrightarrow\quad
  d \approx \frac{0{,}5 - g(\vb{p})}{\norm{\nabla g(\vb{p})}},
\end{equation}
avec $d > 0$ quand la nouvelle surface est à l'extérieur. Loin de la surface ($g$ proche de 0 ou de 1), on prend une
carte de distance exacte au voxel près (transformée de Maurer \cite{maurer2003}).

\begin{table}[htbp]
\centering\small
\caption{Sensibilité au seuil sur le lévrier : écart moyen (en valeur absolue) entre la surface à 226~HU et celle à
176~HU, et 95\textsuperscript{e} centile de l'écart avec la surface à 276~HU.}
\label{tab:sensibilite}
\begin{tabular}{lrr}
\toprule
\titrecol{Os (série)} & \titrecol{écart moyen, 176 HU (mm)} & \titrecol{95\textsuperscript{e} centile, 276 HU (mm)} \\
\midrule
fémur (hanche--genou) & 0,10 & 0,38 \\
tibia (genou--cheville) & 0,11 & 0,40 \\
fibula (genou--cheville) & 0,17 & 0,40 \\
rotule (genou--cheville) & 0,20 & 0,46 \\
talus (genou--cheville) & 0,21 & 0,39 \\
\bottomrule
\end{tabular}
\end{table}

Les grands os bougent en moyenne d'un dixième de millimètre ; 95~\% des points restent à moins de 0,5~mm. Les petits
os fins (fibula, rotule, talus) sont un peu plus sensibles.

\section{Les rapports}

Pour chaque série, l'étape 5 écrit dans \fichier{resultats/<série>/} :
\begin{itemize}
  \item \fichier{rapport.md} (lisible) et \fichier{rapport.json} (toutes les valeurs) : la série, les os exportés
  (volume, triangles, étanchéité, Euler, affinage, sensibilité), l'identification avec ses raisons, les avertissements ;
  \item \fichier{apercu_3d.png} et \fichier{apercu_coupes.png}, pour un contrôle visuel rapide ;
  \item \fichier{reperes_auto.mrk.json} : la tête du fémur, au format de 3D Slicer.
\end{itemize}

\begin{danslecode}
\fichier{code_source/etape4_maillage.py} : \code{remplir()}, \code{mailler()} (flou, Flying Edges, lissage,
décimation, normales), \code{affiner_demi_hauteur()}, \code{controle()} (étanchéité, Euler, volume),
\code{bords_touches()}, \code{ecrire_stl()}, \code{ecrire_seg_nrrd()} et \code{lire_seg_nrrd()}, \code{exporter()}.
\fichier{code_source/etape5_controle.py} : \code{sensibilite_seuil()}, \code{apercu_3d()}, \code{apercu_coupes()},
\code{rapport_md()}, \code{controler()}.
\end{danslecode}

\begin{attention}[title={Le rapport recopie le nom du patient}]
\fichier{rapport.md} et \fichier{rapport.json} recopient le nom et l'identifiant lus dans l'en-tête DICOM. C'est
utile pour reconnaître un examen, mais ces fichiers ne doivent pas être partagés tels quels (\cref{sec:securite}).
\end{attention}
